% ============================================================
\documentclass[master.tex]{subfiles}

\begin{document}

\ifSubfilesClassLoaded{%
  \begin{center}
    {\LARGE\bfseries\color{isublue}
     Chapter 15:
     Dynamic Treatment Regimes
     }\\[6pt]
    {\large Theory and Methods in Causal Inference}\\[4pt]
    {\normalsize Department of Statistics, Iowa State University}
  \end{center}
  \bigskip
  \hrule height 1.2pt \relax
  \smallskip
}{%
  \chapter{Dynamic Treatment Regimes}
}

\section{Why Dynamic Treatment Regimes?}
\label{w15:sec:motivation}

Chapters~\ref{ch:policy-eval} and~\ref{ch:optimal-policy} studied treatment rules in a single-stage setting, where treatment is assigned once on the basis of baseline covariates. In many scientific and clinical applications, however, treatment decisions are made repeatedly over time. Later decisions may depend on intermediate outcomes, adherence, toxicity, or other evolving patient characteristics \citep{moodie2007}.

A \emph{dynamic treatment regime} (DTR) is a sequence of decision rules that assigns treatment at each stage as a function of the information available up to that stage. The goal is no longer to evaluate or learn a single one-time treatment decision, but to study a policy that adapts over time.

This chapter has two parts. First, we define and identify the value of a fixed dynamic regime using the potential-outcomes framework and the longitudinal g-formula. Second, we study how to learn an optimal regime by exploiting the recursive structure of multistage decision problems.

The key new difficulty is that future treatment decisions are themselves part of the intervention. As a result, both the counterfactual notation and the optimization problem become inherently sequential.

\subsection*{Learning Objectives}
By the end of this chapter, students should be able to:
\begin{enumerate}
    \item define a dynamic treatment regime as a sequence of treatment rules indexed by patient history;
    \item define counterfactual outcomes and regime value in multistage settings;
    \item state the assumptions required for identification in longitudinal causal inference;
    \item derive the longitudinal g-formula for a fixed regime;
    \item explain regression-based, weighting-based, and doubly robust estimation of regime value;
    \item explain backward induction and dynamic programming in multistage treatment optimization;
    \item describe multistage Q-learning and A-learning;
    \item explain value search estimation over parametric regime classes and its backward iterative implementation;
    \item describe the classification analogy for optimal regime estimation with binary treatment;
    \item describe the role of SMART designs in developing adaptive interventions;
    \item recognize the main inferential and practical challenges in dynamic treatment regime analysis.
\end{enumerate}

\section{Worked Example: Responder/Nonresponder Re-randomization}
\label{w15:sec:example}

To see why dynamic treatment regimes are needed, consider a two-stage clinical setting. At baseline, each patient is randomized to an initial treatment. After some period, the patient's intermediate response is assessed. Responders may continue their current therapy, whereas nonresponders may be re-randomized to an alternative second-stage treatment.

In this setting, the treatment decision at stage 2 depends on information revealed after stage 1. Thus the treatment strategy is not a single action but a sequence of rules: a first-stage rule that chooses initial treatment from baseline covariates, and a second-stage rule that chooses follow-up treatment based on baseline covariates, initial treatment, and intermediate response.

This example illustrates three features that will drive the rest of the chapter:
\begin{enumerate}
  \item treatment is assigned at multiple times;
  \item later decisions depend on observed history;
  \item the quality of an early decision depends on what can be done later.
\end{enumerate}
The third point is especially important. The best initial treatment cannot be chosen myopically; it must be evaluated in light of the future treatment options that become available after observing intermediate outcomes.

To make this concrete, suppose a clinician is treating patients with depression. At baseline, each patient has covariates \(X_1\), such as age, symptom severity, and prior treatment history. At the first decision point, the clinician chooses an initial treatment
\[
A_1 \in \{0,1\},
\]
where \(A_1=0\) denotes standard psychotherapy and \(A_1=1\) denotes medication.

After several weeks, the patient is assessed for response. Let
\[
X_2
\]
denote the intermediate clinical status, which may include a response indicator, side effects, and updated symptom scores. At the second decision point, treatment may be modified based on this intermediate information. For example,
\[
A_2 \in \{0,1\},
\]
where \(A_2=0\) denotes continuation of the current treatment and \(A_2=1\) denotes augmentation or switching.

The final outcome \(Y\) might be a depression score measured at the end of follow-up, where larger values indicate better clinical improvement.

A dynamic treatment regime in this setting consists of two decision rules:
\[
d=(d_1,d_2),
\]
where
\[
d_1(X_1)
\]
specifies the initial treatment based on baseline covariates, and
\[
d_2(H_2)
\]
specifies the second-stage treatment as a function of the full stage-2 history $H_2$, which is formally defined in Section~\ref{w15:sec:setup} and includes the baseline covariates, the initial treatment, and the intermediate response.

For example, consider the following regime:
\begin{quote}
``Start with psychotherapy. If the patient responds well, continue the same treatment; if the patient does not respond, switch to medication.''
\end{quote}
This regime may be written as
\[
d_1(X_1)=0,
\]
and
\[
d_2(H_2)=
\begin{cases}
0, & \text{if the patient responds at stage 2,}\\
1, & \text{if the patient does not respond.}
\end{cases}
\]

The regime $d = (d_1, d_2)$ captures features~1 and~2 above directly. Feature~3 is the deeper point: in a two-stage setting, the value of $d_1$ depends on what $d_2$ will do with the intermediate outcome that $A_1$ helps to produce. The best initial treatment therefore cannot be identified by looking at stage~1 alone. This is why optimal dynamic regimes require backward induction rather than one-shot optimization; the formal recursion is developed in Section~\ref{w15:sec:dp}.

\section{Histories, Actions, and Regimes}
\label{w15:sec:setup}

We now formalize the multistage setting introduced in Section~\ref{w15:sec:example}. Suppose treatment is assigned at stages $k = 1,\ldots,K$. Let $A_k \in \mathcal{A}_k$ denote the treatment at stage $k$, and let $X_k \in \mathcal{X}_k$ denote the covariates or intermediate outcomes observed before treatment at stage $k$. In particular, $X_1$ contains baseline covariates available before any treatment is administered. The final outcome $Y$ is measured after Decision $K$.

For subject $i$, the observed data consist of
\[
O_i = (X_{1i},A_{1i},X_{2i},A_{2i},\ldots,X_{Ki},A_{Ki},Y_i).
\]
In the depression example of Section~\ref{w15:sec:example}, $K = 2$: $X_1$ is the vector of baseline clinical characteristics, $A_1 \in \{0,1\}$ is the initial treatment, $X_2$ is the intermediate response status, $A_2 \in \{0,1\}$ is the second-stage treatment decision, and $Y$ is the final depression score.

\subsection{Overbar notation}

Sequences of stage-indexed variables arise throughout the analysis of dynamic treatment regimes. We adopt the following compact notation, standard in the literature \citep{robins1986new}.

For $k = 1,\ldots,K$, define the cumulative sequences and their sample spaces by
\[
\bar{x}_k = (x_1,\ldots,x_k), \quad \bar{a}_k = (a_1,\ldots,a_k),
\qquad
\bar{\mathcal{X}}_k = \mathcal{X}_1\times\cdots\times\mathcal{X}_k, \quad \bar{\mathcal{A}}_k = \mathcal{A}_1\times\cdots\times\mathcal{A}_k.
\]
By convention, $\bar{x} = \bar{x}_K$ and $\bar{a} = \bar{a}_K$ denote the complete histories. For a regime $d = (d_1,\ldots,d_K)$, write $\bar{d}_k = (d_1,\ldots,d_k)$.

In the two-stage setting, $\bar{x}_2 = (x_1, x_2)$ and $\bar{a}_1 = a_1$ are the only nontrivial cumulative sequences; the complete histories are $\bar{x} = (x_1, x_2)$ and $\bar{a} = (a_1, a_2)$.

\subsection{Histories}

\begin{defn}{Observed history}
\label{w15:defn:history}
The observed history available at Decision $k$ is
\[
h_1 = x_1 \in \mathcal{H}_1 = \mathcal{X}_1,
\qquad
h_k = (\bar{x}_k, \bar{a}_{k-1}) \in \mathcal{H}_k = \bar{\mathcal{X}}_k \times \bar{\mathcal{A}}_{k-1},
\quad k = 2,\ldots,K.
\]
\end{defn}

The history $h_k$ records all covariate and treatment information accrued through Decision $k$. In the two-stage setting, $H_1 = X_1$ and $H_2 = (X_1, A_1, X_2)$.

\subsection{Dynamic treatment regimes}

\begin{defn}{Dynamic treatment regime}
\label{w15:defn:dtr}
A $K$-stage dynamic treatment regime is a sequence of decision rules
\[
d = (d_1,\ldots,d_K), \qquad d_k : \mathcal{H}_k \to \mathcal{A}_k, \quad k = 1,\ldots,K.
\]
That is, $d_k(h_k) = d_k(\bar{x}_k, \bar{a}_{k-1}) \in \mathcal{A}_k$ specifies the treatment to be assigned at Decision $k$ given history $h_k$.
\end{defn}

A static regime assigns the same treatment sequence to every subject, whereas a dynamic regime allows treatment to depend on evolving patient history. The defining feature of a dynamic regime is that each rule $d_k$ uses only information available up to stage $k$: it maps the observed history $H_k$ into a treatment decision without access to future outcomes or future covariate values.

\begin{remark}
The recursive structure of histories implies that applying regime $d$ forward in time requires substituting earlier rules into later ones. For instance, with $K = 3$, the treatment selected at Decision 3 is
\[
d_3\bigl\{\bar{x}_3,\; d_1(x_1),\; d_2\bigl(\bar{x}_2, d_1(x_1)\bigr)\bigr\},
\]
since the treatments at Decisions 1 and 2 become part of the history for Decision 3. This recursive dependence is precisely why multi-stage optimization requires backward induction; see Section~\ref{w15:sec:dp}.
\end{remark}

\subsection{Feasible action sets and regime classes}

In many medical applications, not all treatment options in $\mathcal{A}_k$ are clinically available to every patient at Decision $k$. For example, in the depression setting of Section~\ref{w15:sec:example}, augmenting treatment makes sense only for patients who have not yet responded. More generally, the set of options that are medically or ethically permissible may depend on the patient's accrued history.

\begin{defn}{Feasible set}
\label{w15:defn:feasible}
For $k = 1,\ldots,K$, the \emph{feasible set} at Decision $k$ given history $h_k \in \mathcal{H}_k$ is a subset
\[
\Psi_k(h_k) \subseteq \mathcal{A}_k.
\]
The collection $\Psi = (\Psi_1,\ldots,\Psi_K)$ encodes the history-dependent treatment constraints across all stages.
\end{defn}

A regime $d = (d_1,\ldots,d_K)$ is called \emph{$\Psi$-specific} if $d_k(h_k) \in \Psi_k(h_k)$ for every $h_k$ in the support of the history distribution, and for every $k = 1,\ldots,K$. The relevant class of regimes is then
\[
\mathcal{D} = \{d = (d_1,\ldots,d_K) : d_k(h_k) \in \Psi_k(h_k) \text{ for all } k \text{ and all } h_k\}.
\]
The optimal regime is defined as $d^{\mathrm{opt}} \in \mathcal{D}$ satisfying $V(d^{\mathrm{opt}}) \geq V(d)$ for all $d \in \mathcal{D}$.

In practice, one often restricts attention to a smaller class of candidate regimes — for example, linear threshold rules $d_k(h_k) = \mathbf{1}(\eta_k^\top h_k > 0)$, tree-based rules, or other interpretable decision classes. The optimal regime is then defined relative to that restricted class. The choice of regime class involves a bias--variance tradeoff: richer classes can represent more complex optimal boundaries, but are harder to estimate from finite samples.

\begin{remark}
Feasible sets interact with the positivity assumption: the requirement $P\{A_k = a_k \mid H_k = h_k\} > 0$ need only hold for actions $a_k \in \Psi_k(h_k)$, not for the entire action space $\mathcal{A}_k$. When $\Psi_k(h_k) = \mathcal{A}_k$ for all $k$ and $h_k$, the feasibility constraint is vacuous and recovers the setup of Definition~\ref{w15:defn:dtr}.
\end{remark}

\section{Potential Outcomes under a Regime}
\label{w15:sec:counterfactuals}

In the single-stage setting (Chapters~\ref{ch:policy-eval}--\ref{ch:optimal-policy}), potential outcomes are indexed by a single treatment level, such as $Y^*(1)$ or $Y^*(0)$. In a multistage setting, treatment decisions occur sequentially, so the counterfactual outcome must be indexed by an entire regime. The potential outcomes are accordingly \emph{nested}: the stage-$(k+1)$ variables are defined under the treatments assigned at stages $1,\ldots,k$, and the final outcome is defined under the full sequence of induced decisions.

\subsection{Nested potential outcomes}

For a fixed treatment sequence $\bar{a} = (a_1,\ldots,a_K) \in \bar{\mathcal{A}}$, the potential outcomes at each stage are:
\begin{itemize}
  \item $X_k^*(\bar{a}_{k-1})$: the intermediate covariates that would be observed between Decisions $k-1$ and $k$ if treatments $\bar{a}_{k-1} = (a_1,\ldots,a_{k-1})$ had been administered, for $k = 2,\ldots,K$;
  \item $Y^*(\bar{a})$: the final outcome that would be observed if the complete treatment sequence $\bar{a}$ had been administered.
\end{itemize}
In the two-stage setting these reduce to $X_2^*(a_1)$ and $Y^*(a_1,a_2)$.

Collectively, we denote the full set of potential outcomes for all treatment sequences by
\[
W^* = \bigl\{X_k^*(\bar{a}_{k-1}),\, Y^*(\bar{a});\; \bar{a}_k \in \bar{\mathcal{A}}_k,\; k = 1,\ldots,K\bigr\}.
\]
This notation will be used to state the sequential ignorability assumption compactly in Section~\ref{w15:sec:identification}.

In the two-stage case, one may think of applying a regime $d = (d_1, d_2)$ as proceeding in four steps:
\begin{enumerate}
  \item apply the stage-1 rule $d_1$ to baseline history $H_1$;
  \item observe the resulting counterfactual intermediate response $X_2^*(d_1(H_1))$;
  \item apply the stage-2 rule $d_2$ to the counterfactual history $H_2^*(d_1)$;
  \item observe the corresponding counterfactual outcome $Y^*(d_1(H_1),\, d_2(H_2^*(d_1)))$.
\end{enumerate}
This nesting is not merely a notational complication. It reflects the causal fact that future treatment decisions are themselves part of the intervention: to evaluate a regime one must account for how the intermediate outcomes produced by earlier decisions feed into later decisions.

\subsection{Counterfactual outcome under a regime}

\begin{defn}{Counterfactual outcome under a regime}
\label{w15:defn:yd}
For a $K$-stage regime $d = (d_1,\ldots,d_K) \in \mathcal{D}$, the counterfactual outcome under regime $d$ is
\[
Y^*(d) = Y^*\!\bigl(d_1(H_1),\, d_2(H_2^*(d_1)),\, \ldots,\, d_K(H_K^*(\bar{d}_{K-1}))\bigr),
\]
where $H_k^*(\bar{d}_{k-1}) = (\bar{X}_k^*(\bar{d}_{k-1}),\, \bar{d}_{k-1})$ is the stage-$k$ history that would arise if regime $d$ were followed through Decision $k-1$. For the two-stage case this simplifies to $Y^*(d) = Y^*(d_1(H_1), d_2(H_2^*(d_1)))$.
\end{defn}

\begin{defn}{Regime value}
\label{w15:defn:regime-value}
The value of a dynamic treatment regime $d \in \mathcal{D}$ is
\[
V(d) = \E\{Y^*(d)\}.
\]
The optimal regime $d^{\mathrm{opt}} \in \mathcal{D}$ is any regime satisfying $V(d^{\mathrm{opt}}) \geq V(d)$ for all $d \in \mathcal{D}$.
\end{defn}

The target of the first half of this chapter is the regime value $V(d)$: the mean outcome that would be obtained if regime $d$ were implemented in the target population. Sections~\ref{w15:sec:identification}--\ref{w15:sec:value-estimation} establish when and how $V(d)$ can be identified and estimated from observed data for a \emph{fixed} regime $d$. The second half of the chapter then turns to finding the regime that maximizes $V(d)$.

\section{Identification in Longitudinal Settings}
\label{w15:sec:identification}

We now ask when the value of a fixed dynamic regime $d$ can be identified from observed longitudinal data. The required assumptions are the multistage analogues of the standard causal assumptions from the single-stage setting, and were formalized by \citet{robins1986new} under the name \emph{sequential randomization assumptions} (SRA).

\emph{Sequential consistency} requires that if the observed treatment sequence coincides with what regime $d$ prescribes, then the observed intermediate variables and final outcome equal the corresponding counterfactual quantities under $d$. \emph{Sequential ignorability} requires that at each stage $k$, conditional on the observed history $H_k$, the treatment assignment $A_k$ is independent of all future counterfactual variables. This is the longitudinal analogue of no unmeasured confounding. \emph{Sequential positivity} requires that at each stage $k$, every action prescribed by the regime must occur with positive probability among units with the corresponding observed history.

\subsection{Sequential consistency}

\begin{defn}{Sequential consistency}
\label{w15:defn:seq-consistency}
The observed intermediate covariates and final outcome coincide with the corresponding potential outcomes under the treatment sequence actually received:
\[
X_k = X_k^*(\bar{A}_{k-1}), \quad k = 2,\ldots,K,
\qquad
Y = Y^*(\bar{A}).
\]
\end{defn}

Sequential consistency requires that there is a single well-defined version of each potential outcome, so that the label $\bar{a}$ fully specifies the intervention. In the two-stage setting this reduces to $X_2 = X_2^*(A_1)$ and $Y = Y^*(A_1,A_2)$.

\subsection{Sequential ignorability}

\begin{defn}{Sequential ignorability}
\label{w15:defn:seq-ignorability}
For $k = 1,\ldots,K$,
\[
W^* \indep A_k \mid \bar{X}_k, \bar{A}_{k-1},
\]
where $W^*$ denotes the full collection of potential outcomes, and $\bar{A}_0$ is vacuous.
\end{defn}

Sequential ignorability is the longitudinal analogue of no unmeasured confounding: at each Decision $k$, the treatment assigned is conditionally independent of all potential outcomes given the history $H_k = (\bar{X}_k, \bar{A}_{k-1})$ accrued up to that point. The assumption implies that every variable used to inform the treatment decision at stage $k$ is measured and included in $H_k$.

In the two-stage setting, sequential ignorability implies
\[
\{Y^*(a_1,a_2), X_2^*(a_1)\} \indep A_1 \mid H_1,
\qquad
Y^*(a_1,a_2) \indep A_2 \mid H_2,\; A_1 = a_1,
\]
for all $(a_1,a_2) \in \mathcal{A}_1 \times \mathcal{A}_2$. The compact $W^*$ formulation makes clear that the same single assumption governs all potential outcomes simultaneously.

\begin{warn}{Sequential Ignorability Is Untestable}
Sequential ignorability is not testable from the observed data. In observational studies it must be defended on substantive grounds. The assumption is automatically satisfied in a SMART design (Section~\ref{w15:sec:smart}), because treatment is randomized at each stage conditional on the observed history.
\end{warn}

\subsection{Sequential positivity}

\begin{defn}{Sequential positivity}
\label{w15:defn:seq-positivity}
For $k = 1,\ldots,K$,
\[
P\{A_k = a_k \mid H_k = h_k\} > 0
\]
for every $(h_k, a_k)$ with $h_k$ in the support of the history distribution and $a_k \in \Psi_k(h_k)$.
\end{defn}

Sequential positivity requires that every feasible treatment option has a positive probability of being observed at every reachable history. When the feasible sets $\Psi_k(h_k)$ are strict subsets of $\mathcal{A}_k$, the condition applies only to options within $\Psi_k(h_k)$.

\begin{remark}
Longitudinal causal inference is harder than single-stage causal inference because the observed history $H_k$ includes intermediate variables $X_2,\ldots,X_k$ that are simultaneously consequences of earlier treatments and predictors of future treatment. Naive regression adjustment that conditions on such variables can introduce collider bias; the g-formula and IPW approaches in Section~\ref{w15:sec:value-estimation} are specifically designed to avoid this pitfall.
\end{remark}

Together, the three assumptions identify $V(d)$ for any fixed regime $d \in \mathcal{D}$. They do not yet tell us how to find the optimal regime; that optimization problem is taken up in Sections~\ref{w15:sec:dp}--\ref{w15:sec:alearning} through backward induction. The three assumptions play distinct roles: sequential consistency links observed and counterfactual data; sequential ignorability removes hidden stage-specific confounding given the observed history; and sequential positivity ensures that the observed data provide support for the regime being evaluated.

\section{The Longitudinal g-Formula}
\label{w15:sec:gformula}

Under sequential consistency, sequential ignorability, and sequential positivity, the value of any fixed regime $d \in \mathcal{D}$ can be expressed entirely in terms of observed-data quantities. This identifying functional, called the \emph{longitudinal g-formula}, was first derived by \citet{robins1986new}. It is the multistage generalization of the standardization formula encountered in single-stage causal inference.

The conceptual logic of the formula is the same in both settings. To evaluate the value of a fixed regime $d$, one must:
\begin{enumerate}
  \item fix the treatment choices at each stage according to the rules in $d$;
  \item propagate the induced distribution of intermediate covariates under those treatment choices;
  \item average the final conditional outcome under the resulting longitudinal law.
\end{enumerate}
The formula is more elaborate than in the single-stage case because the intervention affects not only the final outcome directly, but also the future covariate histories over which we must integrate.

\subsection{Two-stage case}

To build intuition we first state the result for $K=2$.

\begin{thm}{Longitudinal g-formula, $K=2$}
\label{w15:thm:gformula2}
For a fixed two-stage regime $d=(d_1,d_2)$, under sequential consistency, sequential ignorability, and sequential positivity,
\[
V(d)
=
\E\!\left[
  \E\!\left\{
    \E[Y\mid H_2,\,A_2=d_2(H_2)]
  \mid H_1,\,A_1=d_1(H_1)
  \right\}
\right].
\]
\end{thm}

\begin{proof}[Proof sketch]
By iterated expectation and sequential consistency, $V(d)=\E\{Y^*(d)\}=\E\bigl[\E\{Y^*(d)\mid H_1\}\bigr]$. Conditioning on the stage-2 history that would arise under $d_1$ and applying sequential ignorability at stage 2 identifies the inner conditional mean by the observed-data quantity $\E(Y\mid H_2, A_2 = d_2(H_2))$. A second application of sequential ignorability at stage 1 then identifies the outer conditional expectation, yielding the nested standardization formula.
\end{proof}

The three-step logic is visible in the nested structure: the innermost expectation $\E[Y \mid H_2, A_2 = d_2(H_2)]$ fixes stage-2 treatment according to $d$ (step~1); the middle expectation propagates the distribution of $H_2$ under stage-1 treatment $d_1(H_1)$ (step~2); and the outer expectation averages over the baseline distribution of $H_1$ (step~3).

\subsection{General \texorpdfstring{$K$}{K}-stage case}

The two-stage result extends naturally to $K$ stages, where the g-formula takes the form of a product of conditional densities integrated over the covariate paths generated by $d$.

\begin{thm}{Longitudinal g-formula, general $K$}
\label{w15:thm:gformula}
For a fixed $K$-stage regime $d \in \mathcal{D}$, under sequential consistency, sequential ignorability, and sequential positivity,
\[
V(d)
=
\int_{\bar{\mathcal{X}}}
\E\!\bigl\{Y \mid \bar{X}=\bar{x},\, \bar{A}=\bar{d}(\bar{x})\bigr\}
\prod_{k=2}^{K}
p_{X_k\mid\bar{X}_{k-1},\bar{A}_{k-1}}\!\bigl\{x_k \mid \bar{x}_{k-1},\, \bar{d}_{k-1}(\bar{x}_{k-1})\bigr\}
\, p_{X_1}(x_1)\,d\mu(\bar{x}),
\]
where $\mu$ is a dominating measure on $\bar{\mathcal{X}}$ and $\bar{d}(\bar{x}) = \bigl(d_1(x_1), d_2(\bar{x}_2, d_1(x_1)), \ldots\bigr)$ denotes the treatment sequence prescribed by $d$ along the covariate path $\bar{x}$.
\end{thm}

\begin{proof}[Proof sketch]
Under sequential ignorability at stage $k$, the conditional density of $X_k^*(\bar{a}_{k-1})$ given past potential outcomes equals the observed-data conditional density of $X_k$ given $(\bar{X}_{k-1}, \bar{A}_{k-1}) = (\bar{x}_{k-1}, \bar{a}_{k-1})$; see \citet{robins1986new}. Applying this substitution at each stage and integrating over covariate paths yields the product formula above. The two-stage formula of Theorem~\ref{w15:thm:gformula2} is the special case $K=2$.
\end{proof}

\begin{remark}
The longitudinal g-formula is the sequential analogue of standardization. Each factor $p_{X_k|\bar{X}_{k-1},\bar{A}_{k-1}}\{x_k \mid \bar{x}_{k-1}, \bar{d}_{k-1}(\bar{x}_{k-1})\}$ describes how intermediate covariates would evolve if the treatments prescribed by $d$ were administered, and the outer integral averages over the resulting covariate paths. Students should view this as longitudinal standardization: we are averaging conditional mean outcomes over the history distribution induced by the regime, just as single-stage standardization averages over the marginal covariate distribution. The product structure in the formula arises precisely because the history distribution changes at each stage in response to the regime's treatment assignments.
\end{remark}

\section{Estimation of Regime Value}
\label{w15:sec:value-estimation}

We now turn from identification to estimation of the value of a fixed dynamic regime $d$. The main estimators are multistage analogues of the familiar methods from single-stage policy evaluation: outcome regression (g-computation), inverse probability weighting, and augmented inverse probability weighting. The new challenge is that the sequential structure can amplify instability and misspecification across stages, so the practical performance of each estimator depends on how errors accumulate through the recursion.

\subsection{Outcome regression and g-computation}
\label{w15:subsec:gcomp}

The most direct estimator replaces the conditional expectations appearing in the g-formula by estimated regression models and then simulates or averages recursively under the candidate regime. This approach, known as \emph{g-computation} \citep{robins1986new}, mirrors the identifying formula directly. However, it can be sensitive to misspecification of the stage-specific regression models, and such errors can compound through the sequential simulation as each stage builds on earlier fitted models.

Specifically, specify parametric models
\[
p_{X_k|\bar{X}_{k-1},\bar{A}_{k-1}}(x_k \mid \bar{x}_{k-1}, \bar{a}_{k-1};\, \xi_k), \quad k = 2,\ldots,K,
\qquad
\E\{Y \mid \bar{X}=\bar{x}, \bar{A}=\bar{a};\, \xi_{K+1}\}.
\]
Estimate each $\xi_k$ by maximizing the corresponding partial likelihood from the observed data. Then, for a given regime $d$, approximate $V(d)$ by the following Monte Carlo algorithm:

\begin{enumerate}
  \item Draw $x_{1r}$ from the empirical distribution of $X_1$ (or from a fitted model $p_{X_1}(x_1;\hat\xi_1)$).
  \item Draw $x_{2r}$ from $p_{X_2|X_1,A_1}\{x_2 \mid x_{1r},\, d_1(x_{1r});\, \hat\xi_2\}$.
  \item For $k = 3,\ldots,K$, draw $x_{kr}$ from $p_{X_k|\bar{X}_{k-1},\bar{A}_{k-1}}\{x_k \mid \bar{x}_{k-1,r},\, \bar{d}_{k-1}(\bar{x}_{k-1,r});\, \hat\xi_k\}$.
  \item Compute $\hat{y}_r = \E\{Y \mid \bar{X}=\bar{x}_r, \bar{A}=\bar{d}(\bar{x}_r);\,\hat\xi_{K+1}\}$.
  \item Repeat steps 1--4 for $r = 1,\ldots,M$ and set $\widehat{V}_{\mathrm{g}}(d) = M^{-1}\sum_{r=1}^M \hat{y}_r$.
\end{enumerate}

G-computation is consistent when all component models are correctly specified, but it can be computationally intensive and is sensitive to model misspecification in any one of the $K+1$ component models.

\subsection{Inverse probability weighting}
\label{w15:subsec:ipw}

An alternative approach uses inverse probability weighting. The key idea is to reweight observed trajectories so that units whose observed treatment path agrees with the regime represent other similar units who would have followed that path under the target intervention. In multistage settings, the weight is a product of stage-specific inverse treatment probabilities. This product structure is a major practical challenge: even moderate weight instability at each individual stage can accumulate rapidly across stages, leading to high-variance estimators.

Formally, define the \emph{regime consistency indicator}
\[
\delta_i(d)
= \prod_{k=1}^K \mathbb{I}\{A_{ki} = d_k(H_{ki})\},
\]
which equals 1 if and only if subject $i$ received the complete treatment sequence prescribed by $d$. The corresponding product propensity is
\[
\pi_d(\bar{X})
= \prod_{k=1}^K \pi_{d,k}(H_k),
\qquad
\pi_{d,k}(H_k) = P\{A_k = d_k(H_k) \mid H_k\}.
\]
Under sequential consistency, sequential ignorability, and sequential positivity, the IPW estimator
\[
\widehat{V}_{\mathrm{IPW}}(d)
= \frac{1}{n}\sum_{i=1}^n
\frac{\delta_i(d)}{\hat\pi_d(\bar{X}_i)}\,Y_i
= \frac{1}{n}\sum_{i=1}^n
\frac{\delta_i(d)}{\prod_{k=1}^K \hat\pi_{d,k}(H_{ki})}\,Y_i
\]
is consistent when each stage-$k$ propensity model $\hat\pi_{d,k}$ is correctly specified.

\begin{warn}{Variance Inflation from Product Weights}
Product weights grow increasingly variable as $K$ increases. As the number of stages grows, few subjects will have $\delta_i(d) = 1$, and the weights for those subjects can be extreme. This variance inflation is a principal practical limitation of the IPW approach in multistage settings; see Section~\ref{w15:sec:challenges}.
\end{warn}

\subsection{Augmented inverse probability weighting}
\label{w15:subsec:aipw}

Augmented inverse probability weighting combines outcome modeling and weighting, just as in the single-stage setting. Its goal is to retain consistency if either the outcome model or the treatment model is correctly specified at each stage, while improving efficiency relative to pure IPW. In the multistage setting this augmentation is especially valuable because pure product weights can be highly unstable and pure g-computation can be sensitive to recursive misspecification.

The missing-data structure of the problem clarifies the form the augmentation takes. A subject who followed the regime through Decision $k-1$ but then deviated at Decision $k$ has revealed their potential intermediate covariates $X_2^*(d_1),\ldots,X_k^*(d_{k-1})$ but not their potential outcome $Y^*(d)$. This is precisely the structure of monotone missing data \citep{robins1994estimation}: subject $i$ ``drops out'' of the regime at the first stage where $A_{ki} \neq d_k(H_{ki})$.

This missing-data framing motivates the \emph{augmented IPW} (AIPW) estimator, which incorporates partial information from all subjects. For general $K$, the AIPW estimator takes the form
\[
\widehat{V}_{\mathrm{AIPW}}(d)
= \widehat{V}_{\mathrm{IPW}}(d)
+ \frac{1}{n}\sum_{i=1}^n \sum_{k=1}^K
\left(1 - \frac{\delta_{ki}}{\hat\pi_{d,k}(H_{ki})}\right)
\frac{\prod_{j=1}^{k-1}\delta_{ji}}{\prod_{j=1}^{k-1}\hat\pi_{d,j}(H_{ji})}
L_k(H_{ki}),
\]
where $\delta_{ki} = \mathbb{I}\{A_{ki} = d_k(H_{ki})\}$ and $L_k(H_k)$ is an arbitrary function at stage $k$. The choice $L_k \equiv 0$ recovers the plain IPW estimator.

The efficient choice of augmentation function is
\[
L_k(H_k) = \E\{Y^*(d) \mid H_k^*(\bar{d}_{k-1}) = h_k\},
\]
the expected counterfactual outcome under $d$ given the history through stage $k$. In the two-stage case this simplifies to
\begin{align*}
\widehat{V}_{\mathrm{AIPW}}(d)
&= \frac{1}{n}\sum_{i=1}^n \widehat{Q}_{d,1}(H_{1i})\\
&\quad + \frac{1}{n}\sum_{i=1}^n
\frac{\delta_{1i}}{\hat\pi_{d,1}(H_{1i})}
\Bigl\{\widehat{Q}_{d,2}(H_{2i}) - \widehat{Q}_{d,1}(H_{1i})\Bigr\}\\
&\quad + \frac{1}{n}\sum_{i=1}^n
\frac{\delta_{1i}\delta_{2i}}{\hat\pi_{d,1}(H_{1i})\hat\pi_{d,2}(H_{2i})}
\Bigl\{Y_i - \widehat{Q}_{d,2}(H_{2i})\Bigr\},
\end{align*}
where $\widehat{Q}_{d,k}(h_k)$ is a fitted model for $\E\{Y^*(d) \mid H_k^*(\bar{d}_{k-1})=h_k\}$.

\begin{thm}{Double robustness of the AIPW estimator}
\label{w15:thm:dr}
$\widehat{V}_{\mathrm{AIPW}}(d)$ is consistent for $V(d)$ if, at each stage $k$, either the propensity model $\hat\pi_{d,k}$ or the outcome model $\widehat{Q}_{d,k}$ is correctly specified. It is consistent and semiparametrically efficient when both sets of models are correctly specified.
\end{thm}

\begin{remark}
The three-term structure of the two-stage AIPW estimator has a transparent interpretation. The first term is a g-computation estimator anchored at stage 1. The second term corrects for covariate imbalance between the regime-consistent and regime-inconsistent subgroups at stage 2. The third term corrects for any remaining discrepancy between the observed outcome and the stage-2 outcome model for subjects who followed the regime through both stages.
\end{remark}

\subsection{Marginal structural models}

Marginal structural models provide a useful working-model framework for summarizing the effect of dynamic regimes when the regime class is parameterized. Rather than estimating the value of one specific regime, they model the value $V(d_\eta)$ as a function of the regime parameter $\eta$, typically fitting the model by solving weighted estimating equations that use stabilized inverse probability weights. They do not replace the potential-outcomes identification framework; rather, they provide a convenient modeling layer for causal effects of treatment histories under longitudinal weighting schemes. They are especially useful when the scientific goal is to summarize population-level regime effects rather than to identify fully individualized decision rules.

\medskip
The next section turns from fixed-regime estimation to finding the optimal regime.

\section{Dynamic Programming and Backward Induction}
\label{w15:sec:dp}

We now move from evaluating fixed regimes to learning an optimal regime. This requires a new idea: in multistage problems, optimal treatment decisions must be chosen recursively.

The central insight is that the best action at an early stage depends on what can be done optimally at later stages. Therefore, the optimization problem cannot be solved myopically in forward time, stage by stage. Instead, one must work backward from the final stage, treating the optimal future value as the effective outcome when making the current decision.

\subsection{Second-stage decision}

In a two-stage setting, consider stage 2 first. For each possible history $H_2$, define the stage-2 Q-function
\[
Q_2(H_2,A_2)=\E\{Y\mid H_2,A_2\}.
\]
Under the identification assumptions of Section~\ref{w15:sec:identification}, this observed-data conditional expectation equals the causal quantity $\E\{Y^*(A_1,A_2)\mid H_2, A_2\}$, so optimizing over $A_2$ in $Q_2$ yields the optimal causal action.
Then the optimal second-stage action is
\[
d_2^\ast(H_2)=\arg\max_{a_2\in\mathcal A_2} Q_2(H_2,a_2).
\]

Define the optimal continuation value
\[
V_2(H_2)=\max_{a_2\in\mathcal A_2} Q_2(H_2,a_2).
\]

\subsection{First-stage decision}

Now define the first-stage Q-function by
\[
Q_1(H_1,A_1)=\E\{V_2(H_2)\mid H_1,A_1\}.
\]
Then the optimal first-stage action is
\[
d_1^\ast(H_1)=\arg\max_{a_1\in\mathcal A_1} Q_1(H_1,a_1).
\]

Thus backward induction solves the multistage optimization problem by converting it into a sequence of single-stage optimization problems, each using the optimal future value as its outcome. This recursive structure is the basis of multistage Q-learning, A-learning, and value search methods, all of which are developed in the sections that follow. From this point onward, the chapter is concerned primarily with optimal-regime learning, not fixed-regime evaluation.

\section{Q-Learning for Multistage Regimes}
\label{w15:sec:qlearning}

Backward induction gives the conceptual solution to the optimal-regime problem. Q-learning turns that idea into an estimable procedure by fitting a sequence of stage-specific regression models.

The basic strategy is to begin at the final stage, where the outcome $Y$ is directly observed, and then work backward through earlier stages using the estimated optimal future value as a pseudo-outcome. In each step the same tool is used: regression of the current-stage target on the current-stage history and treatment.

\subsection{Stage 2}

At the final stage there is no further decision after treatment $A_2$. Thus the stage-2 problem is structurally the same as a single-stage policy-learning problem, conditional on history $H_2$. Specify a working model
\[
Q_2(H_2,A_2;\beta_2)\approx \E\{Y\mid H_2,A_2\},
\]
estimate $\hat\beta_2$ by regression of $Y$ on $(H_2, A_2)$, and define the estimated optimal second-stage rule by choosing the action with the larger fitted value:
\[
\hat d_2(H_2)=\arg\max_{a_2\in\mathcal A_2} Q_2(H_2,a_2;\hat\beta_2).
\]

The corresponding estimated optimal continuation value is
\[
\widehat V_2(H_2)=\max_{a_2\in\mathcal A_2} Q_2(H_2,a_2;\hat\beta_2).
\]

\subsection{Stage 1}

At stage 1, the final outcome $Y$ is no longer the correct regression target, because the value of the first-stage action depends on what can be done optimally at stage 2. Instead, the target becomes the estimated continuation value $\widehat V_2(H_2)$.

Treat $\widehat V_2(H_2)$ as a pseudo-outcome and fit a working model
\[
Q_1(H_1,A_1;\beta_1)\approx \E\{\widehat V_2(H_2)\mid H_1,A_1\}.
\]
The optimal first-stage rule is then
\[
\hat d_1(H_1)=\arg\max_{a_1\in\mathcal A_1} Q_1(H_1,a_1;\hat\beta_1).
\]

In this way, multistage Q-learning reduces the dynamic optimization problem to a sequence of regression problems solved backward in time.

\begin{defn}{Two-stage Q-learning regime}
\label{w15:defn:qlearn-regime}
The two-stage Q-learning regime is
\[
\hat d_Q=(\hat d_1,\hat d_2),
\]
where \(\hat d_2\) is estimated first and \(\hat d_1\) is estimated using the stage-2 optimal-value pseudo-outcome.
\end{defn}

\begin{warn}{Error Propagation in Q-Learning}
A practical weakness of Q-learning is that model misspecification at a later stage propagates backward through the pseudo-outcomes and can distort earlier-stage decisions. If the stage-2 working model for $Q_2$ is misspecified, the pseudo-outcome $\widehat{V}_2(H_2)$ is biased, and regressing on a biased pseudo-outcome at stage 1 will produce a biased $\hat{d}_1$. The recursive structure that makes Q-learning computationally tractable is therefore also the source of its main statistical vulnerability. This phenomenon is illustrated in the lab (Section~\ref{w15:sec:lab}).
\end{warn}

\section{Value Search and the Classification Analogy}
\label{w15:sec:valuesearch}

Q-learning estimates stage-specific outcome regressions and then derives the implied optimal regime. An alternative is to search directly over a class of candidate regimes and choose the one with the largest estimated value.

This approach is called \emph{value search}. It focuses on the decision rule itself rather than on accurate estimation of the full conditional mean structure. The backward-induction idea remains essential, but the optimization is expressed in terms of estimated regime value rather than estimated Q-functions.

\subsection{Parametric regime classes}

In practice, one restricts attention to a class of regimes indexed by a finite-dimensional parameter $\eta = (\eta_1,\ldots,\eta_K)$, for example linear threshold rules $d_k(h_k;\eta_k) = \mathbf{1}(\eta_k^\top h_k > 0)$. Formally,
\[
\mathcal{D}_\eta = \bigl\{d_\eta = (d_1(\cdot\,;\eta_1),\ldots,d_K(\cdot\,;\eta_K))\bigr\},
\]
and value search estimates the optimal parameter by
\[
\hat\eta^{\mathrm{opt}} = \arg\max_{\eta}\; \widehat{V}(d_\eta),
\]
where $\widehat{V}(d_\eta)$ is any consistent estimator of $V(d_\eta)$ --- for instance, the IPW or AIPW estimator of Section~\ref{w15:sec:value-estimation} evaluated at $d = d_\eta$. This introduces a familiar approximation-estimation tradeoff: a richer regime class may better approximate the true optimal regime, but it also increases the difficulty of the optimization and the risk of overfitting.

\subsection{Backward iterative implementation}

Direct joint maximization over all $(\eta_1,\ldots,\eta_K)$ simultaneously is generally computationally infeasible. In multistage settings, value search still respects the recursive structure of the problem: one optimizes later-stage rules first and then uses the resulting continuation values when optimizing earlier-stage rules. At the final stage $K$, the optimization reduces to a single-stage IPW problem with history $H_K$:
\[
\widehat{V}^{(K)}(d_{\eta,K})
= \frac{1}{n}\sum_{i=1}^n
\frac{\mathbb{I}\{A_{Ki} = d_K(H_{Ki};\eta_K)\}}{\hat\pi_{d,K}(H_{Ki})}\,Y_i.
\]
Maximizing over $\eta_K$ yields $\hat\eta_K^{\mathrm{opt}}$. At each earlier stage $k < K$, with downstream rules already fixed, one maximizes an analogous stage-$k$ value estimator over $\eta_k$ alone. Thus even when value is optimized directly, the dynamic nature of the problem still leads naturally to backward recursion.

\begin{remark}
Q-learning is the special case of backward iterative value search in which the stage-$k$ value estimator is the regression-based pseudo-outcome estimator.  Specifically, $\hat{V}^{(k)}(d_{\eta,k}) = n^{-1}\sum_i \hat{Q}_k(H_{ki}, d_k(H_{ki};\eta_k))$, and maximizing over $\eta_k$ recovers the Q-learning rule $\hat{d}_k(H_k) = \arg\max_{a_k} \hat{Q}_k(H_k, a_k)$.  The three main backward iterative approaches correspond to different stage-$k$ objectives: \citet{zhao2015b} developed the IPW-based approach, \citet{zhang2018} developed the AIPW-based approach, and \citet{zhang2018b} developed the regression-based approach with a focus on interpretable linear decision rules.
\end{remark}

\subsection{The classification analogy for binary treatment}
\label{w15:subsec:classification}

For binary treatment $\mathcal{A}_k = \{0,1\}$, the stage-$k$ value-search optimization reduces to a weighted classification problem. Since the optimal action depends only on the sign of the estimated treatment contrast
\[
\hat{C}_{ki} = \widehat{G}_{ki}(1) - \widehat{G}_{ki}(0),
\]
maximizing the stage-$k$ value estimator over $\eta_k$ is equivalent to minimizing the weighted misclassification rate
\[
\frac{1}{n}\sum_{i=1}^n
|\hat{C}_{ki}|\,
\mathbb{I}\bigl\{\mathbb{I}\{\hat{C}_{ki}>0\} \neq d_k(H_{ki};\eta_k)\bigr\}.
\]
This is a weighted classification problem with pseudo-label $\mathbb{I}\{\hat{C}_{ki}>0\}$ and sample weight $|\hat{C}_{ki}|$, developed by \citet{zhao2015b} for IPW-based contrasts and extended to doubly robust contrasts by \citet{zhang2018}. Any classifier can therefore be plugged in to estimate the optimal stage-$k$ rule. This analogy is conceptually useful and computationally convenient, but it should be viewed as a reformulation of the same dynamic decision problem rather than as a separate causal estimand.

\begin{remark}
The key advantage of the classification framing is flexibility: nonparametric classifiers can estimate the optimal regime even when the decision boundary is not linear or smooth.  The key statistical challenge is that the pseudo-labels $\mathbb{I}\{\hat{C}_{ki}>0\}$ and weights $|\hat{C}_{ki}|$ are themselves estimated, so standard classification theory does not directly apply and bespoke analyses of the excess risk are needed.
\end{remark}

\section{A-Learning and Blip Functions}
\label{w15:sec:alearning}

As in the single-stage setting, Q-learning models the full conditional mean structure, whereas A-learning targets only the treatment contrast that is relevant for decision making. In multistage problems, this contrast is represented by a sequence of stage-specific \emph{blip functions}. Students should view multistage A-learning as the recursive extension of the contrast-based methods from Chapter~\ref{ch:optimal-policy}.

A-learning provides a more targeted alternative to Q-learning when the primary goal is to learn the decision rule rather than to model the full outcome surface. The stage-specific blip functions are linked through the backward recursion, however, and cannot be estimated in isolation: the stage-1 blip depends on the de-blipped pseudo-outcome constructed from the stage-2 blip estimate.

\subsection{Stage-2 blip function}

At the final stage, define the stage-2 blip function as the conditional effect of treatment given history $H_2$. For binary treatment this is simply
\[
\gamma_2(H_2) = Q_2(H_2,1) - Q_2(H_2,0),
\]
the estimated advantage of $A_2 = 1$ over $A_2 = 0$ within each history stratum. More generally, one parameterizes the blip as the treatment-relevant part of the Q-function via the decomposition
\[
Q_2(H_2,A_2)=\nu_2(H_2)+\gamma_2(H_2,A_2;\psi_2),
\]
where $\nu_2(H_2)$ is the history-specific baseline and $\gamma_2$ captures the treatment contrast. In the binary case, $\gamma_2(H_2) = \gamma_2(H_2,1;\psi_2) - \gamma_2(H_2,0;\psi_2)$, so both formulations identify the same optimal rule.

The optimal second-stage rule depends only on the sign of the blip:
\[
d_2^*(H_2) = \mathbf{1}\{\gamma_2(H_2) > 0\}
\]
in the binary-treatment case, so correct specification of $\nu_2$ is not required to recover the optimal rule.

\subsection{Stage-1 blip function}

After estimating $\hat\psi_2$, one constructs a stage-1 pseudo-outcome
by removing the estimated stage-2 treatment effect from the observed
outcome---a step called \emph{de-blipping}:
\[
\tilde Y_i
= Y_i - \bigl(A_{2i} - \hat d_2(H_{2i})\bigr)\,
  \hat\gamma_2\bigl(H_{2i}, 1;\hat\psi_2\bigr),
\]
where $\hat\gamma_2(H_2,1;\hat\psi_2)$ is the estimated blip evaluated at
$A_2=1$.  For binary $A_2$, this adjusts each unit's outcome to reflect
what would have been observed under $\hat d_2(H_{2i})$: nothing changes
for units who already received the optimal action; the blip is subtracted
for units who received $A_2=1$ when $\hat d_2(H_{2i})=0$; and the blip
is added back for units who received $A_2=0$ when $\hat d_2(H_{2i})=1$.
The pseudo-outcome $\tilde Y_i$ represents, under the estimated stage-2
rule, what the outcome would have been had every patient received
$\hat d_2(H_{2i})$ at stage 2.

One then decomposes the stage-1 Q-function similarly:
\[
Q_1(H_1,A_1)=\nu_1(H_1)+\gamma_1(H_1,A_1;\psi_1),
\]
and estimates $\psi_1$ from the pseudo-outcomes $\tilde Y_i$ using the
same propensity-centered estimating equations applied at stage 2.

Thus A-learning proceeds backward through a sequence of stage-specific contrast functions, just as Q-learning proceeds backward through a sequence of Q-functions. The difference is conceptual: Q-learning models full conditional means at each stage, whereas A-learning focuses on the treatment contrasts that determine optimal decisions and avoids specifying the baseline functions $\nu_k$.

\section{SMART Designs}
\label{w15:sec:smart}

Sequential Multiple Assignment Randomized Trials (SMARTs) are experimental designs specifically built for studying adaptive treatment strategies. In a SMART, treatment is randomized at multiple decision points according to a pre-specified protocol, with later randomizations often depending on intermediate response status.

\begin{defn}{SMART}
\label{w15:defn:smart}
A SMART is a multistage randomized design in which treatment is randomized at more than one stage, with later randomizations often stratified by intermediate response status.
\end{defn}

SMART designs are especially important for this chapter because they provide a setting in which the sequential treatment mechanism is known by design. Unlike observational longitudinal data, where sequential ignorability must be defended on substantive grounds, SMART data satisfy the identifying assumptions of Section~\ref{w15:sec:identification} automatically at every decision node.

\subsection{Responder/nonresponder re-randomization}

A prototypical SMART randomizes all participants at stage 1, assesses intermediate response, and then re-randomizes selected subgroups at stage 2. Responders may be re-randomized between continuing the current treatment and augmenting it, while nonresponders may be re-randomized between switching and intensifying. This creates data that are naturally aligned with clinically meaningful adaptive strategies. For example, in a SMART for attention-deficit/hyperactivity disorder (ADHD) in children \citep{murphy2005experimental}, participants are first randomized between medication and behavioral therapy; after eight weeks, responders continue or augment their therapy while nonresponders switch or intensify, generating data on several adaptive strategies simultaneously.

\subsection{Embedded regimes}

A single SMART typically contains several \emph{embedded dynamic treatment regimes}, each corresponding to one possible path through the staged randomization protocol. A regime $d = (d_1, d_2)$ is embedded if, for every history $h_k$ that can arise under the protocol, the treatment $d_k(h_k)$ is among the options randomized at that node. Data from one trial therefore inform multiple candidate adaptive strategies simultaneously. In the ADHD example, one embedded regime is ``give medication; if nonresponding, switch to behavioral therapy'' and another is ``give behavioral therapy; if nonresponding, intensify''; both are estimable from a single trial.

\subsection{Design and analysis}

The major analytical advantage of SMART data is that the stage-specific randomization probabilities $\pi_{d,k}(H_k)$ are known from the design. Sequential ignorability therefore follows from the trial protocol rather than from observational adjustment arguments, and the IPW estimator of Section~\ref{w15:subsec:ipw} can be evaluated at known rather than estimated propensity scores. This strengthens the double-robustness guarantee of the AIPW estimator: any residual bias must come from the outcome model alone, with no propensity estimation error. Similarly, Q-learning regressions at each stage are free of unmeasured confounding, and the centered-treatment estimating equations of A-learning are unbiased without any propensity model.

This does not eliminate all statistical difficulties. Sparsity in certain history strata, variance inflation from product weights, and model misspecification in the outcome regressions remain live concerns. However, the design removes the most fundamental identification concern, making SMART data the cleanest setting for learning dynamic treatment regimes from finite samples. Sample-size planning must additionally account for the prevalence of response at the interim assessment, since this determines the effective sample size for second-stage comparisons.

\section{Practical and Inferential Challenges}
\label{w15:sec:challenges}

Dynamic treatment regimes are appealing because they formalize individualized sequential decision making. They are also substantially more difficult than the single-stage setting, both causally and statistically. This section collects the main practical and inferential challenges that arise in multistage analysis. The challenges fall into three broad groups: identification difficulty, estimation instability, and inference difficulty.

\medskip
\noindent\textbf{A.\enspace Identification difficulty}
\medskip

\subsection{Time-varying confounding}

In observational longitudinal data, a variable measured between two treatment stages may simultaneously be a consequence of the earlier treatment and a predictor of the later treatment. Conditioning on such a variable in a naive regression can induce collider bias, opening spurious pathways rather than closing confounding ones. This is the hallmark of time-varying confounding and the reason the g-formula and IPW approaches of Section~\ref{w15:sec:value-estimation} are needed in place of standard regression adjustment.

\subsection{Positivity deterioration}

Sequential positivity requires every feasible treatment to occur with positive probability at every reachable history. Even when each stage satisfies a reasonable positivity condition marginally, the probability of observing an entire treatment path under a specific regime can be very small: positivity conditions multiply across stages, so support for complete treatment paths decays rapidly with $K$. In practice this leaves some regimes weakly or not at all supported in the observed data, making value estimation unreliable or infeasible for those regimes.

\medskip
\noindent\textbf{B.\enspace Estimation instability}
\medskip

\subsection{Variance inflation}

Inverse probability weights in the multistage setting take the form of a product of $K$ stage-specific propensity factors. Even moderate weight variability at each individual stage compounds across stages, so the resulting product weights can be extremely variable. This is not merely a technical nuisance: in finite samples, a small number of units with very large weights can dominate the IPW estimator, producing estimates with poor mean squared error even when the estimator is nominally consistent. The lab of Section~\ref{w15:sec:lab} illustrates this phenomenon quantitatively.

\subsection{Model misspecification}

In regression-based methods such as Q-learning, a misspecified working model at stage $k$ produces a biased pseudo-outcome $\widehat{V}_k$, which is then used as the regression target at stage $k-1$. The misspecification therefore propagates backward through the entire estimation chain. Thus the multistage structure amplifies the impact of model error relative to the single-stage case, where a misspecified outcome model affects only the final estimate directly.

\medskip
\noindent\textbf{C.\enspace Inference difficulty}
\medskip

\subsection{Nonregularity}

Dynamic regime estimation can produce nonregular asymptotics. At histories where two treatment options have nearly equal Q-function values, the estimated optimal action is unstable: small perturbations in the data can flip the estimated decision. The estimator of the optimal regime is not differentiable in the underlying parameters at these boundary points, so standard smooth asymptotic approximations fail and confidence intervals based on the normal approximation may have poor coverage. Inferential methods robust to this nonregularity, such as the $m$-out-of-$n$ bootstrap or adaptive confidence intervals, are an active area of research.

\begin{remark}
Although dynamic treatment regimes share vocabulary with reinforcement learning, the present chapter remains within the statistical causal inference framework. Counterfactuals, identification assumptions, and semiparametric estimation are central. The distinction is that DTR analysis begins from identification conditions on observational or experimental data rather than from a known transition model.
\end{remark}

\section{Lab: A Two-Stage Dynamic Treatment Regime}
\label{w15:sec:lab}

This lab uses a two-stage simulation to illustrate both parts of the chapter:
\begin{enumerate}
  \item estimation of the value of a fixed dynamic regime;
  \item learning of an optimal regime by backward induction.
\end{enumerate}
The goal is not only to compare methods numerically, but also to expose the practical challenges introduced by sequential decision making. Complete R code is provided in the accompanying file \texttt{lab\_chapter15.R}.

\subsection{Data-generating process}
\label{w15:sec:lab:dgp}

The data are generated as follows for each of $n$ independent units:
\begin{align*}
  X_1     &\sim N(0,1), \\
  A_1     &\sim \operatorname{Bernoulli}(0.5), \\
  X_2\mid A_1 &\sim N(-A_1,\, 1), \\
  A_2     &\sim \operatorname{Bernoulli}(0.5), \\
  Y       &= 2A_1 X_1 + 2A_2 X_2 + \varepsilon, \quad \varepsilon\sim N(0,1).
\end{align*}
Both treatment stages are randomized ($A_1, A_2$ each balanced at 0.5),
so sequential ignorability holds by design.
The key structural feature is that $A_1 = 1$ \emph{lowers} $X_2$ by one
unit on average, creating a negative indirect effect of stage-1 treatment
on the final outcome through the intermediate covariate.

This DGP is designed to support the two-layer analysis of the lab. For fixed-regime evaluation (Part~A), any specific regime $d = (d_1, d_2)$ can be evaluated by g-computation, IPW, or AIPW against the known oracle value. For optimal-regime learning (Part~B), the true optimal regime is analytically tractable, providing exact benchmarks against which Q-learning and A-learning can be assessed.

\paragraph{True optimal regime.}
The stage-2 Q-function is $Q_2(H_2, A_2) = 2A_1 X_1 + 2A_2 X_2$, so the
stage-2 blip is $\gamma_2(H_2, A_2) = 2A_2 X_2$.  The optimal stage-2 rule
is
\[
  d_2^*(H_2) = \mathbf{1}(X_2 > 0).
\]
The pseudo-outcome under optimal stage-2 treatment is $V_2^*(H_2)
= 2A_1 X_1 + 2\max(0, X_2)$.  Taking its conditional expectation over
$X_2\mid A_1 \sim N(-A_1, 1)$, the total stage-1 blip is
\[
  Q_1(H_1,1) - Q_1(H_1,0) = 2X_1
    + 2\bigl[\mathbb{E}[\max(0,X_2)\mid A_1=1]
             - \mathbb{E}[\max(0,X_2)\mid A_1=0]\bigr]
    = 2X_1 - 0.632,
\]
where the constant $-0.632$ is the indirect cost of $A_1 = 1$ through
the reduced distribution of $X_2$.  The optimal stage-1 rule is therefore
\[
  d_1^*(X_1) = \mathbf{1}(X_1 > 0.316).
\]

\paragraph{Reference values.}
\begin{align*}
  V(d^*)        &= 1.32, \\
  V(\text{never treat}) &= 0, \\
  V(\text{always treat}) &= -2.00.
\end{align*}
The large regret of ``always treat'' reflects the fact that $A_1 = 1$
depresses $X_2$, making stage-2 treatment ineffective or harmful for most
patients.

\subsection{Candidate methods}
\label{w15:sec:lab:methods}

\paragraph{Part A: Fixed-regime evaluation.}
For a fixed target regime $d$ (e.g., the true optimal $d^* = (d_1^*, d_2^*)$), we compare three estimators of $V(d)$:
\begin{enumerate}
  \item \textbf{G-computation}: parametric outcome regression models $\E\{Y \mid H_2, A_2\}$ and $\E\{X_2 \mid A_1\}$ are fitted, and $V(d)$ is estimated by Monte Carlo simulation under the regime.
  \item \textbf{IPW}: the product-weight estimator $\widehat{V}_{\mathrm{IPW}}(d)$ of Section~\ref{w15:subsec:ipw}, with known propensity scores $\pi_{d,k} = 0.5$.
  \item \textbf{AIPW}: the doubly robust estimator $\widehat{V}_{\mathrm{AIPW}}(d)$ of Section~\ref{w15:subsec:aipw}, combining the fitted outcome models and known propensity scores.
\end{enumerate}
This comparison directly illustrates product-weight instability (Q4 below) and the efficiency gain from augmentation.

\paragraph{Part B: Optimal-regime learning.}
For optimal-regime learning, we compare five strategies evaluated by oracle value on a fresh Monte Carlo sample:
\begin{enumerate}
  \item \textbf{Optimal oracle}: the true optimal regime
        $(d_1^*, d_2^*)$ evaluated on the true model; this is an
        upper bound, attainable only with full knowledge of the DGP.
  \item \textbf{Always treat} / \textbf{never treat}: two static regimes
        that assign $A_k \equiv 1$ or $A_k \equiv 0$ at every stage.
  \item \textbf{Q-learning}: two-stage backward induction as described in
        Section~\ref{w15:sec:qlearning}, with working linear models that
        include the $A_k \times X_k$ interaction term at each stage.
  \item \textbf{A-learning}: two-stage blip estimation as described in
        Section~\ref{w15:sec:alearning}, with blip models
        $\gamma_k(H_k, A_k;\psi_k) = A_k(\psi_{k0} + \psi_{k1}X_k)$
        and propensity-centered estimating equations.
\end{enumerate}

\subsection{Simulation results}
\label{w15:sec:lab:results}

The results are organized to match the two-layer structure of the lab. Part~A examines fixed-regime value estimation; Part~B examines optimal-regime learning. All results are based on 400 Monte Carlo replications.

\paragraph{Part A: Fixed-regime estimation.}
For fixed-regime evaluation of $d^*$, g-computation and AIPW both achieve low bias at $n = 500$, with AIPW showing modestly lower variance due to its doubly robust construction. The IPW estimator is unbiased but has substantially higher variance, driven by product-weight instability documented in Table~\ref{w15:tab:cv} and Q4 below.

\paragraph{Part B: Optimal-regime learning.}
Table~\ref{w15:tab:sim} reports the mean oracle value, its Monte Carlo
standard deviation, and the regret $V(d^*) - V(\hat d)$ across 400
replications with $n = 500$.

\begin{table}[h]
\centering
\caption{Mean oracle value, standard deviation, and regret for each
method ($n = 500$, 400 Monte Carlo replications). Regret is defined as
$V(d^*) - V(\hat d)$ where $V(d^*)=1.32$ is the true optimal value.}
\label{w15:tab:sim}
\begin{tabular}{lccc}
\toprule
Method & Mean value & SD & Regret \\
\midrule
Optimal (oracle)   & $1.3197$ & $0.009$ & $0.001$ \\
Never treat        & $0.0000$ & $0.000$ & $1.320$ \\
Always treat       & $-2.001$ & $0.021$ & $3.321$ \\
Q-learning         & $1.3169$ & $0.010$ & $0.004$ \\
A-learning         & $1.3149$ & $0.011$ & $0.006$ \\
\bottomrule
\end{tabular}
\end{table}

Both Q-learning and A-learning recover nearly all of the optimal value,
with regrets of $0.004$ and $0.006$ respectively.
The two static regimes are substantially suboptimal.
``Always treat'' performs worst because $A_1 = 1$ systematically
reduces $X_2$, lowering the stage-2 payoff regardless of patient
characteristics.

\subsection{Questions for exploration}
\label{w15:sec:lab:questions}

\paragraph{Q1. Why does backward induction produce a better stage-1 decision than naive regression?}
The naive alternative to backward induction regresses $Y$ directly on
$(H_1, A_1)$ at stage 1, without first forming the pseudo-outcome
$\widehat{V}_2(H_2)$.  This naive regression reflects the \emph{average}
future payoff under a balanced random $A_2$, not the \emph{optimal}
future payoff.  As a result, it \emph{overstates} the indirect cost of
$A_1 = 1$ through $X_2$.  In this DGP, with $A_2 \sim \operatorname{Bernoulli}(0.5)$
independent of $(X_1, A_1, X_2)$, the naive stage-1 blip is
\begin{align*}
  \E[Y \mid H_1, A_1=1] - \E[Y \mid H_1, A_1=0]
  &= 2X_1 + 2\,\mathbb{E}[A_2 X_2 \mid A_1=1]
     - 2\,\mathbb{E}[A_2 X_2 \mid A_1=0] \\
  &= 2X_1 + 2 \cdot 0.5 \cdot (-1) - 2 \cdot 0.5 \cdot 0
   = 2X_1 - 1,
\end{align*}
giving a decision threshold of $X_1 > 0.5$.
Backward induction instead uses the pseudo-outcome
$\widehat{V}_2(H_2) = \max(Q_2(H_2,0),Q_2(H_2,1))$, which captures the
value of \emph{optimally} adapting $A_2$ to the observed $X_2$.
Under the optimal rule $d_2^*(H_2) = \mathbf{1}(X_2 > 0)$, the stage-2
treatment partially compensates for the reduction in $X_2$ caused by
$A_1 = 1$, so the indirect cost is correctly estimated as $-0.632$
rather than $-1$, and the threshold is lowered to $0.316$.
A simulation with $n = 5{,}000$ confirms this: the BI threshold is $0.30$
and the naive threshold is $0.48$, close to the theoretical values.
Naive is too conservative---it under-treats patients with $X_1 \in
(0.316,\, 0.5)$ who should be treated---because it does not account for
the fact that optimal stage-2 treatment will adapt to the observed $X_2$.
Across 400 replications with $n = 500$, BI achieves a mean value of
$1.317$ versus $1.303$ for the naive approach, a gain of $0.014$.

\paragraph{Q2. How does the stage-2 continuation value feed back into the stage-1 optimal choice?}
The stage-2 pseudo-outcome $\widehat{V}_2(H_2) = \max(Q_2(H_2, 0),
Q_2(H_2, 1))$ encodes the value of optimal future treatment.  When
$A_1 = 1$, $X_2$ is lower on average, so $\widehat{V}_2$ is also lower.
In the simulation ($n = 5{,}000$), the mean $\widehat{V}_2$ is $0.80$ when
$A_1 = 0$ and only $0.13$ when $A_1 = 1$---a drop of $0.67$.  This cost
is directly absorbed by the stage-1 regression, which learns that the
direct gain $2X_1$ must exceed $0.632$ to justify treating at stage 1.
Without backward induction the indirect cost is invisible to the stage-1
model, and the threshold is set too low.

\paragraph{Q3. How does stage-2 model misspecification propagate backward to corrupt stage-1 Q-learning?}
We refit Q-learning with the $A_2 \times X_2$ interaction omitted from
the stage-2 working model.  The misspecified model estimates the stage-2
blip as approximately constant and negative---reflecting the average
$\mathbb{E}[Y \mid A_2 = 1] - \mathbb{E}[Y \mid A_2 = 0] = \mathbb{E}[X_2] \approx -0.5$
over the data---and consequently sets $\hat{d}_2 \equiv 0$ for every unit.
Across 400 replications with $n = 500$, the mean value drops from $1.317$
(correct model) to $0.702$ (misspecified model), a loss of $0.615$.
This illustrates a key hazard of backward induction: a misspecified
stage-$k$ model corrupts $\widehat{V}_k$, which in turn corrupts the
pseudo-outcome fed to stage $k-1$, so misspecification at any stage
propagates backward through the entire estimation chain.

\paragraph{Q4. How variable are product IPW weights, and why does variance grow with the number of stages?}
With $A_1, A_2 \sim \operatorname{Bernoulli}(0.5)$ and any deterministic
regime $d$, the effective IPW weight for unit $i$ is
\[
  w_i = \frac{\delta_i(d)}{\pi_{d,1}(H_{1i})\,\pi_{d,2}(H_{2i})},
\]
where $\delta_i(d) = \mathbf{1}(A_{1i} = d_1(X_{1i}))\,
\mathbf{1}(A_{2i} = d_2(H_{2i}))$.
Since $\pi_{d,k} \approx 0.5$, unit $i$ receives weight $w_i \approx 4$
if it follows $d$ at both stages (probability $0.25$) and $w_i = 0$
otherwise (probability $0.75$).  The population coefficient of variation
is $\sqrt{3} \approx 1.73$, independently of $n$.
The simulation confirms convergence to this theoretical value as $n$
grows (Table~\ref{w15:tab:cv}), demonstrating that weight instability is
an intrinsic feature of the product structure rather than a
finite-sample artifact.  With $K$ stages the product has $K$ propensity
factors and the CV grows with $K$, making pure IPW estimation increasingly
unreliable in long-horizon regimes; see Section~\ref{w15:sec:challenges}.

\begin{table}[h]
\centering
\caption{Coefficient of variation (CV) of effective product IPW weights
for the optimal regime $d^*$, across 200 replications at each sample size.
Theoretical value: $\sqrt{3} \approx 1.732$.}
\label{w15:tab:cv}
\begin{tabular}{rcc}
\toprule
$n$ & Mean CV & SD of CV \\
\midrule
$100$   & $1.819$ & $0.213$ \\
$250$   & $1.762$ & $0.133$ \\
$500$   & $1.755$ & $0.099$ \\
$1000$  & $1.738$ & $0.063$ \\
$2000$  & $1.736$ & $0.045$ \\
\bottomrule
\end{tabular}
\end{table}

\section{Summary}
\label{w15:sec:summary}

This chapter extended single-stage policy methods to dynamic treatment regimes, where treatment decisions are made sequentially using evolving patient history. The chapter has two parts, corresponding to two distinct problems: evaluating the value of a fixed dynamic regime, and learning the optimal regime.

\medskip
\noindent\textbf{Part I.\enspace Fixed-regime causal inference}
\medskip

\begin{enumerate}

  \item \textbf{Dynamic treatment regimes generalize single-stage rules.}
    A $K$-stage DTR is a sequence of decision rules
    $d = (d_1,\ldots,d_K)$, each mapping the observed history $H_k$ to a
    treatment action. The counterfactual outcome $Y^*(d)$ is nested: each
    stage-$k$ variable is defined under the treatments assigned at stages
    $1,\ldots,k-1$. The regime value $V(d) = \E\{Y^*(d)\}$ is the
    expected outcome if $d$ were implemented in the target population.

  \item \textbf{Identification requires three sequential assumptions.}
    Sequential consistency ($X_k = X_k^*(\bar{A}_{k-1})$, $Y = Y^*(\bar{A})$),
    sequential ignorability ($W^* \indep A_k \mid \bar{X}_k, \bar{A}_{k-1}$
    for all $k$), and sequential positivity together imply the
    \emph{longitudinal g-formula}:
    \[
      V(d) = \int \E\!\bigl\{Y \mid \bar{X}=\bar{x}, \bar{A}=\bar{d}(\bar{x})\bigr\}
             \prod_{k=2}^K p_{X_k\mid\bar{X}_{k-1},\bar{A}_{k-1}}
             \{x_k \mid \bar{x}_{k-1}, \bar{d}_{k-1}(\bar{x}_{k-1})\}\,
             p_{X_1}(x_1)\,d\mu(\bar{x}).
    \]
    This is the multistage analogue of standardization.
    Sequential ignorability is the longitudinal analogue of no unmeasured
    confounding; it holds by design in SMART trials.

  \item \textbf{Regime value can be estimated three ways.}
    G-computation evaluates the g-formula by Monte Carlo simulation under
    parametric models, but is sensitive to misspecification in any of its
    $K+1$ component models.  The IPW estimator reweights the outcomes of
    subjects who followed the regime by the inverse product propensity; it
    is consistent under a correct propensity model but suffers from
    variance inflation as $K$ grows.  The AIPW estimator achieves double
    robustness: it is consistent if, at each stage, either the propensity
    or the outcome model is correctly specified, and is semiparametrically
    efficient when both are correct.

\end{enumerate}

\medskip
\noindent\textbf{Part II.\enspace Optimal-regime learning}
\medskip

\begin{enumerate}
  \setcounter{enumi}{3}

  \item \textbf{Backward induction is the key to optimal regime
    estimation.}  The optimal regime is found by working backward from the
    final stage: the optimal stage-$K$ action maximizes the Q-function
    $Q_K(H_K,A_K) = \E\{Y\mid H_K,A_K\}$; the optimal continuation value
    $V_K(H_K)$ is then used as a pseudo-outcome to define $Q_{K-1}$, and
    so on.  This dynamic-programming recursion ensures that every
    stage-$k$ decision accounts for the value of optimal future actions.

  \item \textbf{Q-learning, A-learning, and value search are three
    approaches to estimation.}  Q-learning fits parametric working models
    for each Q-function and substitutes estimated values into the
    backward-induction recursion \citep{zhang2018b}.  A-learning focuses
    on the treatment contrast (blip function) at each stage and constructs
    a de-blipped pseudo-outcome to propagate the stage-$k$ blip estimate
    backward, avoiding unnecessary outcome modeling.  Value search
    maximizes the estimated regime value over a parametric class
    $\mathcal{D}_\eta$; for binary treatment this optimization reduces to
    a weighted classification problem \citep{zhao2015b,zhang2018}.

\end{enumerate}

\medskip
\noindent\textbf{Design and practical challenges}
\medskip

\begin{enumerate}
  \setcounter{enumi}{5}

  \item \textbf{SMART designs provide the cleanest experimental framework
    for DTR learning.}  In a SMART, treatment is randomized at each
    decision point, often with re-randomization based on intermediate
    response status.  Sequential ignorability holds by design, propensity
    scores are known, and each embedded regime is supported by data from
    the full trial.  Careful power calculations must account for the
    prevalence of response at the interim assessment.

  \item \textbf{Several inferential challenges are unique to the
    multistage setting.}  Time-varying confounding requires g-formula or
    IPW methods rather than naive regression adjustment.  Positivity
    deteriorates as the number of stages grows, leading to extreme product
    weights and support problems.  Model misspecification at any stage
    propagates backward through the recursive estimation chain.  When the
    optimal treatment boundary passes through a region of near-zero
    effect, estimators of the optimal regime are nonregular and standard
    asymptotic theory fails.

\end{enumerate}

\medskip
\noindent\textit{Dynamic treatment regime analysis is best viewed as the longitudinal extension of policy evaluation and policy learning. In both parts, the sequential structure of the decision process must be respected: identification requires sequential analogues of the standard causal assumptions, and optimization requires backward induction rather than myopic stage-by-stage maximization.}

\section*{Problems}

\begin{enumerate}
    \item Define a two-stage dynamic treatment regime and explain the difference between a static regime and a dynamic regime.

    \item In the two-stage setting, define \(X_2^*(a_1)\), \(Y^*(a_1,a_2)\), and \(Y^*(d)\) for a regime \(d=(d_1,d_2)\).

    \item State the assumptions of sequential consistency, sequential ignorability, and sequential positivity in a two-stage setting.

    \item Derive the nested-expectation form of the longitudinal g-formula for a fixed two-stage regime \(d=(d_1,d_2)\).

    \item Write down the longitudinal IPW estimator of \(V(d)\) for a fixed two-stage regime.

    \item Explain why multistage optimal treatment requires backward induction.

    \item Describe the steps of two-stage Q-learning.

    \item Explain how multistage A-learning differs from multistage Q-learning.

    \item What is a SMART design, and why is it useful for adaptive intervention development?

    \item Design a simple simulation study for comparing two-stage Q-learning and IPW-based regime evaluation.
\end{enumerate}

\ifSubfilesClassLoaded{%
  \bibliographystyle{plainnat}
  \bibliography{causal_inference}
}{}

\end{document}